\chapter*{勘误表}
\addcontentsline{toc}{chapter}{勘误表}
\markboth{勘误表}{勘误表}%

本表汇总中译全过程（2026 年 8--9 月）中发现的原书排印疑误，共 314 条，按章分组。正文相应位置均以脚注注明（体例四要素：原书印作什么、疑应为什么、此处径改抑或照排存疑，末缀“——译者注”），本表仅撮其要，供对照原书阅读时参考。处置凡分两类：“径改”指译文已按酌定改排补正；“存疑照排”指证据不足或两读可通，译文一仍原书之旧，仅备案待考。2026 年 9 月外部复审两轮所新增、改写与撤销的条目，均已同步入表。所记页码均指原书印刷页码。

\begin{description}[font=\normalfont,leftmargin=6.5em,itemsep=0.15em]

\subsection*{第~1~章}
\item[p.~7] 注记 1.1.2 称二阶导数在一阶导数的\trKernel{}上\trWellDefined{}——一般并不成立（反例 $g(x,y)=x$、$\Phi(u,v)=(u+v^{2},v)$：同一方向两坐标下取值 $0$ 与 $2$），仅\trCriticalPoint{}（$dg=0$）处\trHessian{}方内禀（存疑照排）。
\item[p.~8] §1.1.b 称\trDegenerate{}的\trQuadraticForm{}“在 $0$ 处的二阶导数就是它自己” $\to$ 按本章约定 $(d^{2}f)_{x}(X,Y)=X\cdot(\widetilde{Y}\cdot f)(x)$ 应为其 $2$ 倍（差因子 $2$；“就是它自己”仅在把 \trHessian{}取作 \trTaylor{}展开二次项系数的约定下成立）（存疑照排）。
\item[p.~10] $v_i(u_1,\dotsc,u_{n-d})$ $\to$ $v_i(u_1,\dotsc,u_d)$（径改）。
\item[p.~10] 分块矩阵右上块 $\Bigl(\sum_{k}\frac{\partial v_k}{\partial u_i}\cdot v_\ell\Bigr)$ 疑漏排因子 $t_k$ $\to$ 径补（径改）。
\item[p.~11] $\sum\varepsilon_i^2$ $\to$ $\frac{1}{2c}\sum\varepsilon_i^2$（径改）。
\item[p.~12] \trMorseLemma{} $n=1$ 情形首式 $f=f(0)+\frac{1}{2}f''(0)x^{2}+\varepsilon(x)x^{2}$ 与所给 $\varepsilon$ 的积分（三阶余项）不自洽（取 $f=x^{2}+x^{3}$ 即见首式不成立）；证明仍可经换元自救（存疑照排）。
\item[p.~20] 练习 8 的 $\widetilde{f}$ 对 $z$ 二次齐次而非不变，未先限制到单位\trSphere{}即降至 $\CP^3(\CC)$ 并不\trWellDefined{}（存疑照排）。
\item[p.~20] 练习 9(2) “show that $F$ is a submersion” 漏印 $D_{1}$ $\to$ 按值域与证明所用者，所证的\trSubmersion{}应是映射 $D_{1}F$（径改）。

\subsection*{第~2~章}
\item[p.~25] 边界公式 $\partial_{\pm}U=\{x\in U\mid Q(x)=\pm\varepsilon,\ \dots\}$ 之 $x\in U$ 与 $U$ 的定义 $-\varepsilon<Q(x)<\varepsilon$ 不容，按字面为空集 $\to$ 按后文用法应理解作 $x\in\overline{U}$（或 $\partial U$）（存疑照排）。
\item[p.~33/35] 定理 2.1.11 证明步骤 (2)/(3) 的收缩目标层级与 $H$ 的定义及分解式均印作 $\alpha+\varepsilon$ $\to$ 按 $F^{\alpha+\varepsilon}=V^{\alpha+\varepsilon}$ 本身、$r_{t}$ 仅在 $\varOmega(a)$ 内非恒同，中间层级应作 $\alpha-\varepsilon$（径改）。
\item[p.~34] 同证明的 $F$ 构造式 $\|x_{+}\|^{2}$ 后无减号，与同页 $F(a)=\alpha-\mu(0)$ 及 $dF$ 系数 $-1-\mu'$、$1-2\mu'$ 均不\trCompatible $\to$ 疑漏一枚负号（径改）。
\item[p.~35] $\varepsilon\le\|x_-\|^2\le\varepsilon\|x_+\|^2$ $\to$ 右端应为 $\varepsilon+\|x_+\|^2$（径改）。
\item[p.~41] 引理 2.2.8 陈述仅设 $\varepsilon>0$ $\to$ 取 $\varepsilon>\max f-\alpha_{j}$ 时修改带为空、断言为假（其证明开篇即改取充分小的 $\varepsilon$），疑漏“充分小”之限定（径改）。
\item[p.~43] \trEmbedding{} $\boldsymbol{\Psi}$ 的靶为开带 $f^{-1}\bigl(\,]\alpha_{j}+\varepsilon,\alpha_{j}+2\varepsilon[\,\bigr)$，而 $z=0$、$z=m$ 两端之像却指定落在其两端面上 $\to$ 按其用法宜写作闭带（存疑照排）。
\item[p.~44] $P\cap\{(0,q,0)\mid q\in Q\}$ $\to$ 疑为 $P\cap(\{0\}\times Q\times[0,m])$ 之误（存疑照排）。
\item[p.~45] $\operatorname{codim}_P W^s_{X'}\cap P$ $\to$ $\operatorname{codim}_P(W^s_{X'}\cap P)$（径改）。
\item[p.~46] 解式 $x_{i}(s)$ 右端 $\int_{0}^{-s}\beta_{i}(-t)\,dt$ $\to$ 换元 $u=-t$ 后应为 $\int_{0}^{-s}\beta_{i}(t)\,dt$（取 $s=-m$ 验之两端不相当）（径改）。
\item[p.~48] $f(\varphi_X^{-s}(x))=s$ 与“$X$ \trIncoming”合读有粗疏之嫌（若读作 $X$ 定义于 $\partial V$ 的双侧邻域则自洽）（存疑照排）。
\item[p.~51] “Mores--Bott” $\to$ “Morse--Bott”（径改）。

\subsection*{第~3~章}
\item[p.~55] “can indeed by compactified” $\to$ “can indeed be compactified”（径改）。
\item[p.~58/59] 邻域基定义印作 $0<i_{0}<\dots<i_{k-1}<i_{k}=q$ 并对 $j\le k$ 设 $\mu_{j}\in\mathcal{L}(c_{i_{j}},c_{i_{j+1}})$——首段缺 $a$、且 $j=k$ 时 $i_{k+1}$ 无定义 $\to$ 疑索引统体错位一位，径改作 $0=i_{0}<i_{1}<\dots<i_{k}=q$、$\mu_{j}\in\mathcal{L}(c_{i_{j-1}},c_{i_{j}})$（径改）。
\item[p.~68] $\bigl(C_k(f),\partial_X\bigr)$ $\to$ $\bigl(C_\star(f),\partial_X\bigr)$（径改）。
\item[p.~68] 称定理 3.2.7“含有一条关于\trOrientation{}的陈述”，而英译本该陈述无“\trOrientation”字样（疑法文 orientée 于英译中脱落）（存疑照排）。
\item[p.~70] $\operatorname{Ind}_{f_0}(b)$ $\to$ $\operatorname{Ind}_{f_1}(b)$（径改）。
\item[p.~71] $g\bigl|_{[-1/3,\,2/3]}$ $\to$ $g\bigl|_{[-1/3,\,1/3]}$（径改）。
\item[p.~71/72/74] 自首构造 $X_{i}-\operatorname{grad}g$ 之后的诸\trInterpolation{}与边界公式均印作 $+\operatorname{grad}g$ $\to$ 按\trAdapted{}条件与\trTrajectory{}方向推演疑均应作 $-\operatorname{grad}g$（本节七处一并径改）；又同页 $C_{\star+1}(f_{0})$ $\to$ 由 $\operatorname{Ind}_{\widetilde{F}}(a,0)=\operatorname{Ind}_{f_{0}}(a)+1$ 应作 $C_{\star-1}(f_{0})$（径改）。
\item[p.~71] 命题 3.4.3 按字面断言任意的 $\cC^{1}$ 小\trPerturbation{}仍是\trAdapted{}的——一般不真（局部反例：$f=x^{2}$、$X=-2x\partial_{x}$ 加紧支小\trPerturbation），\trPerturbation{}须限于 $f$ 的\trAdapted{}\trPseudogradient{}所成空间内（存疑照排）。
\item[p.~72] $K$ 的定义域第二因子 $\bigl[\tfrac{1}{3},\tfrac{4}{3}\bigr]$ $\to$ 据 $t\in[-\tfrac{1}{3},\tfrac{1}{3}]$ 时 $K_{s,t}=F_{s}$ 之要求，疑漏一枚负号，应作 $\bigl[-\tfrac{1}{3},\tfrac{4}{3}\bigr]$（径改）。
\item[p.~74] $[-\tfrac{1}{2},\tfrac{4}{3}]$ $\to$ $[-\tfrac{1}{3},\tfrac{4}{3}]$（径改）。
\item[p.~74] $S\circ\partial_{X_0}+\Phi^G\circ\Phi^F+\Phi^H+\partial_{X_2}\circ S=0$ 与下式 $\Phi^G\circ\Phi^F-\Phi^H=S\circ\partial_{X_0}+\partial_{X_2}\circ S$ 在 $\ZZ$ 系数下不\trCompatible $\to$ 本段矩阵计算须先按 $\ZZ/2$ 系数理解；扩到 $\ZZ$ 须对移位\trComplex{}统一调整\trOrientation{}与\trHomotopy{}符号，非单一改一处符号所能解决（存疑照排）。
\item[p.~77] $F(x,t)=f_t(x)+g(x)$ 之 $g(x)$ $\to$ $g(t)$（径改）。
\item[p.~78] 练习 14 末句印作 co-oriented $\to$ 由所设\trExact{}列给出的乃是 $E\cap F$ 的\trOrientation{}而非其在 $W$ 中的\trCoorientation（有显式反例），疑为笔误（径改）。

\subsection*{第~4~章}
\item[p.~81] “to be nonempty”句中 $\mathcal{L}_{X'}(a',b')$ $\to$ $\mathcal{L}_Y(a',b')$（径改）。
\item[p.~82] 命题 4.2.1 及其证明中\trTensorProduct\trComplex{}的\trDifferential{}均印作 $\partial_X\otimes 1+1\otimes\partial_Y$，第二项不带任何符号 $\to$ 对带整数系数的\trGraded{}\trTensorProduct\trComplex，第二项疑应带符号 $(-1)^i$（Koszul 符号约定），否则 $d^2=0$ 一般不成立（存疑照排）。
\item[p.~84] 注记中 $H_2(\CP^3(\RR);\ZZ)$ $\to$ 按上下文所论均为\trMorseHomology{}的\trTorsion{}比较，疑应 $HM_2(\CP^3(\RR);\ZZ)$（径改）。
\item[p.~87] $P_{\CP^2(\RR)}$ 漏排自变量 $\to$ $P_{\CP^2(\RR)}(t)$（径改）。
\item[p.~90] 命题 4.5.7 照排——按字面缺少\trCompact{}（无边界）前提，实为假：取非\trCompact{}的 $V=S^1\times\RR$ 与 $f(\theta,t)=t$，则 $f$ 无任何\trCriticalPoint{}而 $\pi_{1}(V)\cong\ZZ$ 非平凡（存疑照排）。
\item[p.~93] 命题 4.6.3 对任意\trSubmanifold{}包含（含像落在带边目标之边界上者）声称延拓保\trCriticalPoint $\to$ 与场沿边界\trIncoming{}的条件相抵，严格表述须先设\trEmbedding{}像落在目标内部（或先添\trCollar{}）（存疑照排）。
\item[p.~95] 引理 4.6.6 设 $a\in D^{N}$ 照排——证明所用的\trDiffeomorphism{}保持边界，$a$ 应限于 $D^{N}$ 的内点（与练习 19 的开\trDisk{}口径一致）（存疑照排）。
\item[p.~98] $(\psi\circ(\mathrm{Id},v))_\star$ $\to$ $(\mathrm{Id},(\psi,v))_\star$（径改）。
\item[p.~100] \trConnectingHomomorphism{}式印作 $\partial a=\sum_{b}n_{\widetilde{X}}(a,b)$ $\to$ 右端只是模 2 计数之和、漏链基向量 $b$，应作 $\partial[a]=\bigl[\sum_{b}n_{\widetilde{X}}(a,b)\cdot b\bigr]$（径改）。
\item[p.~101] $HM_\star(\overline{W-\mathcal{V}},\mathcal{V})$ 之第二变量 $\to$ $\partial\mathcal{V}$（径改）。
\item[p.~101] 例 4.7.1 的低维端点断言照排存疑：本例须设 $n\ge 2$；§4.8.b 的 Brouwer 论证在 $n=1$ 失效；“不存在 $S^{n}$ 到其子\trSphere{} $S^{k}$（$k<n$）的收缩”的\trHomology{}论证在 $k=0$ 两端皆误；§4.8.c 中 $\partial$ 为\trIsomorphism{}的范围应作 $1\le k\le n$（存疑照排）。
\item[p.~102] 闭圈径称作\trEmbedding{}照排——两段弧在极小点拼接处一般有折角，$u_{\beta}$ 一般只是分段\trSmooth{}的连续\trInjective{}映射（存疑照排）。
\item[p.~106] 末式左端印作 $\partial[a]$ $\to$ $a$ 只是诸 $\alpha_{i}$ 在南半球中的\trLift{}之和、一般并非\trCycle，应作 $\partial[\alpha]$（径改）。
\item[p.~108] “where $\psi$ in an embedding” $\to$ “where $\psi$ is an embedding”（径改）。
\item[p.~109] 定理 4.8.3 证明末\trTangentSpace{}映射印作 $T(\Gamma,\varphi)\widetilde{X}$ $\to$ $\Gamma$ 定义于 $\CP^{n}(\RR)$ 而 $\widetilde{X}$ 是 $S^{n}$ 上的\trVectorField，按前句\trEmbedding{} $(\Gamma\circ p,\varphi)$ 应作 $T(\Gamma\circ p,\varphi)\widetilde{X}$（径改）。
\item[p.~112] \trCompactification{}定义 $\overline{W}^{\,u}(c)=W^{u}(c)\cup\bigl(\overline{\mathcal{L}}(c,d)\times\bigcup_{d}W^{u}(d)\bigr)$ 第一因子 $d$ 悬空 $\to$ 应把 $\bigcup_{d}$ 置于乘积之外：$W^{u}(c)\sqcup\bigcup_{d}\bigl(\overline{\mathcal{L}}(c,d)\times W^{u}(d)\bigr)$（径改）。
\item[p.~113] $W(x,\lambda,U^0,\mathbf{U}^-,\mathbf{U}^+)$ $\to$ $W(\lambda,x,\dotsc)$，实参顺序依定义处统一（径改）。
\item[p.~114] 例 4.9.4(2) 图例称粗线段“由从 $c$ 到 $d$ 的那两条\trTrajectory{}组成” $\to$ 与上文 $\mathcal{L}(c,d)$ 仅含一条\trTrajectory{}相抵，其成分应是 $W^{u}(d)$ 的两条分支（存疑照排）。
\item[p.~116] 定理 4.9.3 证明称 $F$ “把\trCriticalPoint{}送到\trCellular” $\to$ 与陈述中 $F\colon K_\star\to C_\star$（\trCellular{}\trComplex{}到\trMorseComplex{}的\trIsomorphism）方向相反，宜说成把（生成它的）\trCellular{}送到\trCriticalPoint，反向陈述须经 $F^{-1}$（存疑照排）。
\item[p.~117] “(sends $W^u(c,A)$ onto $W^u(c,B)$ and) can also be chosen…”之 and 被置于插语末尾、句法不整 $\to$ 按文义系修饰\trHomeomorphic{}关系的插语，并入正文译出（径改）。
\item[p.~117] 命题 4.9.7 括注的附加结论照排——“$W^{u}(c,A+\varepsilon)$ 映满 $W^{u}(c,A-\varepsilon)$”在紧与非紧之间不可能成立 $\to$ 应改述为\trClosed\trCellular{} $\overline{W}^{\,u}(c,A+\varepsilon)$ 与 $\overline{W}^{\,u}(c,A-\varepsilon)$ 之间的\trHomeomorphism{}并删保截集要求；$\operatorname{Ind}(d)=0$ 情形 $P$ 一般并非整个\trConnectedComponent（存疑照排）。
\item[p.~117] “(that sends $W^u(c,A+\varepsilon)$ onto $W^u(c,A-\varepsilon)$ and) that is the identity on…”同上（径改）。
\item[p.~119] “as describe in Figure 4.11” $\to$ “as described”（径改）。
\item[p.~120] “形如 $(\lambda,x)=x\in W^{u}(c)\cap W^{s}(d)$ 的点”之 $(\lambda,x)=x$ 悬空等号（点对不能等于其第二分量）$\to$ 疑排印衍文，径删等号照文义译出（径改）。
\item[p.~123] $HM_k\bigl(\CP_2(\RR);\ZZ\bigr)$ 之 $\CP$ 误排下标 $\to$ $HM_k\bigl(\CP^2(\RR);\ZZ\bigr)$（径改）。
\item[p.~123] 练习 23 照排——按字面为假（常映射 $V\to\{p\}$ 即到真子集上的收缩），疑本意变形收缩且须补无边界等前提（存疑照排）。

\subsection*{第二部分引言}
\item[p.~127] 首式 $\sum_{k}HM_{k}(W;\ZZ/2)$ 缺 $\dim$（按\trMorse{}不等式口径各项应取维数），且未申明 $W$ \trCompact{}无边界、\trHamiltonian{}量\trNondegenerate{}两假设（存疑照排）。

\subsection*{第~5~章}
\item[p.~131] 命题 5.2.1 照排——按字面尚须默认两事：$W$ 具正维数（$0$ 维\trCompact{}\trManifold{}上零形式即“既\trNondegenerate{}又\trExact{}”）、$W$ 无边界（带边反例即闭\trDisk）（存疑照排）。
\item[p.~133] “indeed defines a 2-form and that we have shown the following result”之 “and that” 疑为赘笔 $\to$ 径按文意译出（径改）。
\item[p.~135] “the speed (or the impulse) of the particle” $\to$ impulse 系法文 impulsion 之英译，按力学语境径译作\trMomentum（径改）。
\item[p.~137] “any vector extending the vector $Z$” $\to$ “any \emph{vector field} extending…”（径改）。
\item[p.~138] 注记 5.4.6 的谱判据印作“$H$ 的\trHessian{}在 $2\pi\ZZ$ 中无\trEigenvalue” $\to$ 应为 $J_{0}\operatorname{Hess}_{x}(H)$ 的\trSpectrum{}与 $2\pi i\ZZ$ 无交（两方向均有反例）（存疑照排）。
\item[p.~138] 命题 5.4.5 证明的导数式印作 $\frac{d}{dt}T_{x}\psi^{t}(Z)=T_{x}\psi^{t}([X_{H},Z])$ $\to$ 实算应带负号（此正负之差不影响非零性推论）（存疑照排）。
\item[p.~139] 注记 5.4.8 “does not have a periodic solution” $\to$ 线性齐次方程零解恒为其解，应指没有非零\trPeriodic{}解（径改）。
\item[p.~140] Hermitian \trInnerProduct{}分解句称\trSymplecticForm{}为其虚部 $\to$ 按 $\langle u,v\rangle=(u,v)-i\omega(u,v)$ 所定之 $\omega$ 乃其负虚部（存疑照排）。
\item[p.~141] \trCalibrated{}\trComplexStructure{}的\trPolarDecomposition{}构造取号相反：由 $(X,AY)=\omega(X,Y)$ 与 $A=BJ$ 所定之 $J$ 使 $\omega(X,JY)=-(BX,Y)<0$，与“被\trCalibrated”的结论相矛盾 $\to$ 疑应自始取 $J=-B^{-1}A$（$\RR^{2}$ 时 $A=-J_{0}$、构造出 $J=-J_{0}$ 可直见）（存疑照排）。
\item[p.~145] 命题 5.6.2 证明 $(3)+(1)$ 式链中间两处等号各漏排一枚负号（二者相消，首尾不误）$\to$ 径补（径改）。
\item[p.~147] 命题 5.6.6 印作 “two (nonzero) integers” $\to$ 负幂 $(A-\lambda\mathrm{Id})^{r}$ 以其\trInvertible{}性为前提，$r$、$s$ 应为正整数（径改）。
\item[p.~148] “the polar decomposition of $\mathrm{GL}(n;\RR)$” $\to$ $\mathrm{GL}(2n;\RR)$（径改）。
\item[p.~148] \trPolarDecomposition{}映射印作 $A\mapsto US$ $\to$ 靶既为乘积集 $\mathrm{U}(n)\times C_{n}$，应作 $A\mapsto(U,S)$；$C_{n}$ 之“开”仅系相对意义（存疑照排）。
\item[p.~149] “Writing $\lambda=\log\ell$” $\to$ $\ell=\log\lambda$（径改）。

\subsection*{第~6~章}
\item[p.~152] 猜想 6.1.2 未假设 $H$ 为 $1$-\trPeriodic，而按真正\trPeriodic{}解使用“\trPeriodic{}为 1 的解” $\to$ 疑应补设 $H$ 为 $1$-\trPeriodic（或把结论改述为时刻 $1$ 映射 $\varphi^{1}$ 的\trFixedPoint{}个数；注记 6.1.3 的重\trParametrization{}只保\trFixedPoint{}不保此类解）（存疑照排）。
\item[p.~153/154] 命题 6.1.5 条件 $\|dX_{H}\|_{L^{2}}<2\pi$ $\to$ 证明所需者是一致算子范数界 $\sup\|dX_{H}\|_{\mathrm{op}}<2\pi$（即 $X_{H}$ 的\trLipschitz{}常数；紧随其后的“\trLipschitz{}常数小于 $2\pi$”一语恰支持此读法）（存疑照排）。
\item[p.~159] 特别地段印作 “$\pi_{2}(W)=0$, that is, $\scr{L}W$ is simply connected” $\to$ $\pi_{2}(W)=0$ 推不出 $\scr{L}W$ \trSimplyConnected（$\mathbf{T}^{2}$ 反例；\trCappingDisk{}独立性只需 $\omega$ 在 $\pi_{2}$ 上为零），“也就是说”疑应删（存疑照排）。
\item[p.~161] 注记 6.3.4(3) “the field $X_H$ is a fixed point of $T_x\varphi^1$” $\to$ 径把\trVectorField{}称作线性映射的“固定点”，按文意应是 $X_H$ 沿\trLoop{}取值所得的向量 $X_H(x)$ 为 $T_x\varphi^1$ 的不动向量（径改）。
\item[p.~161] 引理 6.3.5 缺前提——$H\equiv 0$ 时一切常值\trLoop{}皆\trCriticalPoint{}且作用值相同，两两几何相异者有不可数多条，按字面不成立 $\to$ 应限于后文所需的有限\trNondegenerate{}情形（存疑照排）。
\item[p.~165] \trEnergy{}定义段的\trTrajectory{}方程印作 $\frac{du}{ds}+X(u)=0$，却又称 $u$ 是\trPseudogradient{}场 $X$ 的\trTrajectory $\to$ 按原符号 $u$ 实为 $-X$ 的\trTrajectory、$E(u)\le 0$，与“\trEnergy{}为正、$a$ 位于 $b$ 之上”相抵，应作 $\frac{du}{ds}-X(u)=0$（径改）。
\item[p.~170] 引理 6.5.10 陈述与证明中 $X_t$／$X_\tau$ 混用（积分变量为 $\tau$，被积\trVectorField{}却多印作 $X_t$，仅末行左端作 $X_\tau$），末行中段积分又缺 $d\tau$ $\to$ 按积分变量为 $\tau$ 统一作 $X_\tau$ 并补 $d\tau$（径改）。
\item[p.~171] $\lambda(\dot{u}_k)-X_t(u_k)$ $\to$ $\lambda(\dot{u}_k-X_t(u_k))$（径改）。
\item[p.~176] 引理 6.6.3 首行 $d(y,x_{0})\le 2\varepsilon$ $\to$ 按证明中的迭代至多得几何级数和 $<2\varepsilon_{0}$，疑应作 $d(y,x_{0})\le 2\varepsilon_{0}$（存疑照排）。
\item[p.~179] $\frac{\partial v}{\partial t}=-J(v)\frac{\partial v}{\partial s}$ $\to$ 两导数次序倒置，应作 $\frac{\partial v}{\partial s}=-J(v)\frac{\partial v}{\partial t}$（径改）。
\item[p.~181] 引理 6.6.5 证明三处并案：仅因 $v$ 是 $J$-\trHolomorphic{}的便称 $v^{\star}\omega$ 为\trSymplecticForm、$f$ 为正函数 $\to$ 有分支点时只可断言 $v^{\star}\omega\ge 0$（即 $f\ge 0$；后续论证不受影响）（存疑照排）；楔积印作 $d\theta\wedge d\rho$，按 $ds\wedge dt=\rho\,d\rho\wedge d\theta$ 疑次序相反（存疑照排）；迭次积分 $\int_0^{2\pi}\bigl(\int_0^r f(\rho,\theta)\,\rho\,d\theta\bigr)d\rho$ 内外\trDifferential{}号倒置 $\to$ 据 $A'(r)$ 行应内层对 $\rho$、外层对 $\theta$（径改）。
\item[p.~182/183] 泡泡例三处从简并案：齐次式 $[x^{2},\alpha y^{2},xy]$ 于 $\alpha=0$、$[x,y]=[0,1]$ 处给出 $[0,0,0]$ 无定义；仿射式 $z\mapsto(z,\alpha/z)$ 的定义域应是 $\CC^{\star}$；“诸像曲线都含 $[0,1,0]$”仅当 $\alpha\neq0$（$\alpha=0$ 的像是 $x$ 轴的\trProjective{}闭包）（存疑照排）。
\item[p.~185] “turns it into a cover of the group $\pi_2(W)/\operatorname{Ker}(\omega\colon\pi_2(W)\to\RR)$”——\trCovering{}的对象通常是空间而非群，疑为简笔（存疑照排）。
\item[p.~187] 定理 6.8.1 把 $\scr{L}W$ 写作 $\cC^{\infty}(S^{1};W)$ $\to$ \trContractible{}\trLoop{}所成空间与全体 $\cC^{\infty}$ \trLoop{}仅在 $\pi_{1}(W)=0$ 时相同，宜理解为临时改用全体\trLoop{}、或全限于\trContractible{}的那一\trConnectedComponent（存疑照排）。

\subsection*{第~7~章}
\item[p.~191] 第一步“若 $Z'$ 是沿 $x$ 的另一个\trSymplecticFrame，则道路 $A(t)$ 与 $A'(t)$ \trHomotopic”照排——沿\trLoop{}任意更换\trSymplecticFrame{}并不保证两道路\trHomotopic（绕行一周给 $A'(t)=R(-2\pi t)A(t)$），应补“$Z'$ 可延拓至所选\trCappingDisk{}（或与原\trFrame{}\trHomotopic）”之条件（存疑照排）。
\item[p.~193] $\mathrm{Sp}(2)$ 分类“$\operatorname{tr}A>2\iff$ 相似于 $\operatorname{diag}(\lambda,\lambda^{-1})$，$\lambda>0$”照排——按“当且仅当”读法须排除 $\lambda=1$（此时矩阵即 $\mathrm{Id}$，迹为 $2$）；同理旋转款须排除 $\theta\in2\pi\ZZ$（存疑照排）。
\item[p.~194] 引理 7.1.5 证明选“既\trSymplectic{}又在复共轭下不变的特征向量基”照排——宜理解为\trAdapted{}于复共轭与辛配对的复谱基（$X_{\overline{\lambda}}=\overline{X_{\lambda}}$ 并按辛配对归置），不能读作实特征向量之基（非实\trEigenvalue{}没有实特征向量）（存疑照排）。
\item[p.~196] $W^{\pm}$ 的取法听任任意（“例如……”）$\to$ 任取则\trMaslovIndex{}失去整数性（$n=1$ 取 $W^{+}=R(\pi/2)$ 即对常值道路得 $\mu=-1/2$），应固定为所展示的标准矩阵（存疑照排）。
\item[p.~198] $\exp(tJS_0(\lambda))$ $\to$ $\exp(tJ_0S(\lambda))$（径改）。
\item[p.~200] $R'R^{-1}+RR'^{-1}=0$ $\to$ $R'R^{-1}+R\bigl(R^{-1}\bigr)'=0$（径改）。
\item[p.~200] 证明中 $e^{i\ell t}$ $\to$ 按所设 $S=\operatorname{diag}(\ell\pi,\ell\pi)$ 应为 $e^{i\ell\pi t}$（漏一因子 $\pi$，与 $\mu(e^{tJ_{0}S})=-\ell$ 相配）（径改）。
\item[p.~201] 同证明的分块构造默认 $n\ge 2$（第二分块矩阵需占两枚辛平面）$\to$ $n=1$ 时仅能实现 $k=0$ 与奇数 $k$（存疑照排）。
\item[p.~203] “a basis $(Y_1,\dotsc,Y_k)$ of $E_{\bar\lambda}$” $\to$ $E_{1/\bar\lambda}$（径改）。
\item[p.~204] §7.3.b 首段由 $Q(iX)=Q(X)$ 径断 $E_{\lambda}^{\pm}$ 在乘 $i$ 下稳定 $\to$ $Q$ 之不变只保证正负号不变，任意取定的正（负）子空间未必复稳定，应取相应 Hermitian 型的正负谱空间（存疑照排）。
\item[p.~208/209] 命题 7.3.10 及证明把 $m(\lambda)/2$ 当作可谈奇偶性的整数 $\to$ $\lambda$ 与 $1/\lambda$ 两簇互反、单簇 $m(\lambda)$ 可为奇数，应合并两簇总\trMultiplicity{}后再谈偶性（$\lambda=-1$ 自反簇原论证有效）（存疑照排）。
\item[p.~209] $\lambda_{t_k+1},\dotsc,\lambda_{r_k}$ $\to$ $\lambda^k_{t_k+1},\dotsc,\lambda^k_{r_k}$（径改）。
\item[p.~211] “$(M+iN)v=\lambda v$，从而 $\lambda$ 是 $U$ 的\trEigenvalue” $\to$ $(M-iN)v=\lambda v$、$\bar\lambda$（径改）。
\item[p.~219] \trGeneralizedEigenspace{} $E_\lambda$、$E_{1/\lambda}$ $\to$ $E_\mu$、$E_{1/\mu}$（径改）。
\item[p.~211] $\rho$ 的构造与 $\det_{\CC}U$ 的归一化符号相反（$n=1$、$A=R_{\theta}$ 时按本节定义得 $\rho(A)=e^{-i\theta}$，与性质 (3) 的 $\rho(A)=\det_{\CC}U=e^{i\theta}$ 相抵），且本证明所印 $J_0$ 的两个特征空间标签互相反写 $\to$ 全文性符号约定矛盾，修复须原作者选定一致约定（存疑照排）。
\item[p.~214] 求和下标 $\lambda\in S^1-\RR$ $\to$ $\lambda\notin S^1-\RR$（径改）。
\item[p.~214] $\mathrm{Sp}(2n)^{-}$ 验证式的 $+\pi$ 排在 $i$ 作用域之外（模长 $e^{\pi}\ne 1$，与右端 $(-1)^{m_{0}/2}$ 不符）$\to$ 应作 $\exp\bigl(i(\sum m(\lambda)\pi/2+\pi)\bigr)$（径改）。
\item[p.~216] “恰有两个正实\trEigenvalue”与其所欲证命题的陈述口径不合（存疑照排）。
\item[p.~216] 引理 7.1.5 证明第 (1) 步以普通特征向量配对 $\omega(X,Y)=1$ $\to$ $\lambda$ 亏损（有 Jordan 块）时普通特征向量未必够用，应允许广义特征向量并一致经 $A$ 写出（存疑照排）。
\item[p.~217] $\overline{\beta(s)\lambda(s)}$ 一句中 $\lambda$ 记号前后不一（同句前文作 $\beta(s)\lambda$ 无 $(s)$，$\lambda(s)$ 未定义）$\to$ 按上下文取 $\beta(s)\lambda(s)$ 一致读法（径改）。

\subsection*{第~8~章}
\item[p.~225] 两处 Sobolev 范数 display 的\trDifferential{}号与积分限的归属彼此错位（外层 $-\infty\to+\infty$ 本属 $s$、内层 $0\to 1$ 本属 $t$，紧贴被积函数的却是 $ds$）$\to$ 按累次积分读法应作 $dt\,ds$ 或交换两层积分号（存疑照排）。
\item[p.~228] $y_i\colon\RR\times S^1\to\RR^{2n}$ $\to$ $\RR$（径改）。
\item[p.~228] 引理 8.2.4 以 $\sup_{\RR\times S^{1}}\bigl(|g|,|\partial g/\partial s|,|\partial g/\partial t|\bigr)$ 为常数，按字面一般不足 $\to$ 所缺为一只依赖 $p$ 的普适常数（乘法连续性本身无误）（存疑照排）。
\item[p.~229] \trAtlas{}之卡 $\Phi_{w}$ 的定义域 $\|Y\|\le r_{0}/K$ 印作闭不等号 $\to$ 卡的模型域应开（或视作对球内部取限制）（存疑照排）。
\item[p.~230] \trDiffeomorphism{} $\Psi_i\colon B_i\to\overline{B}(0,1)$ 之目标 $\to$ 疑应开球 $B(0,1)$（诸 $B_i$ 皆开，紧随其后的 sup 恰取于开球 $B(0,1)$）（存疑照排）。
\item[p.~231] \trConvolution{}局部估计的 $\sup_{K}$ 照印——$x\in K$ 处取用的是 $g_{\ell}-f$ 在 $K-\operatorname{Supp}\chi_{k}$ 上的值 $\to$ 严格应作 $\sup_{K-\operatorname{Supp}\chi_{k}}$（补上后结论不受影响）（存疑照排）。
\item[p.~232] $\beta$ “of class $\cC^\infty$ on $\RR^n$”（疑 $\RR^m$，依据不足）（存疑照排）。
\item[p.~233] $\varepsilon'_{k}$ 分母替换式取 $\min_{j}$ 照印 $\to$ 分母须支配诸 $\|\beta\circ\Psi_{i}\circ\Psi_{j}^{-1}\|$ 之最大者（$\min$ 既压不住最大者又可能除零），疑应 $\max_{j}$（存疑照排）。
\item[p.~236] 注记 8.4.2 称展示式首行诸项“零阶”、余下诸项“一阶”——与该式两行的实际阶数及本注记后文分类恰相反 $\to$ 两称谓互换（径改）。
\item[p.~237] $\operatorname{End}(\RR^n)$ $\to$ $\operatorname{End}(\RR^{2n})$（径改）。
\item[p.~240] $y_j\bigl\langle-\frac{\partial Z_i}{\partial t}+(dX_t)(Z_j),Z_i\bigr\rangle$ $\to$ $\frac{\partial Z_j}{\partial t}$（径改）。
\item[p.~241] 命题 8.4.6 证明内 $A(\sum y_{i}Z_{i})$ 展开式首行右端漏印因子 $y_{i}$（次行即恢复出现）$\to$ 径补（径改）。
\item[p.~242] $\cC^{\infty}_{\varepsilon}(W,H_0)$ $\to$ 本节余处皆作 $\cC^{\infty}_{\varepsilon}(H_0)$，疑多一参数 $W$ 系衍余（存疑照排）。
\item[p.~245] 引理 8.5.2 证明首句 $Y\in\cC^{0}_{\infty}(\RR\times S^{1};\RR^{2n})$（上标 $0$、下标 $\infty$）$\to$ 证明明言 $\lim_{s\to\pm\infty}Y=0$，所需即具有紧\trSupport{}的\trSmooth{}函数，疑为 $\cC^{\infty}_{0}$ 之上下标互换（径改）。
\item[p.~247] “the dense open set $R(u)$ of regular values of $u$” $\to$ regular points（径改）。
\item[p.~253] 定理 8.5.7 集合记号 “$y=f(x)$” $\to$ $y=\scr{F}(x)$（径改）。
\item[p.~254] “the Floer equation for $(H_0+h_0,J)$” $\to$ $(H_0+h,J)$（径改）。
\item[p.~256] 链式第二行第二个括号 $\frac{\partial v}{\partial s}$ $\to$ $\frac{\partial v}{\partial t}$（径改）。
\item[p.~256] 命题 8.6.1 证明末取 $\varphi=\psi_1$ $\to$ $\varphi=\psi_1^{-1}$（由 $v=\psi_t^{-1}u$ 与 $u$ 的\trPeriodic{}性，$\psi_{t+1}=\psi_t\circ\psi_1$，得 $v(s,t+1)=\psi_1^{-1}(v(s,t))$）（径改）。
\item[p.~259] “$v$ is locally constant and $(s_0,t_0)\in C(v)$” $\to$ 与上下文相抵，疑应“每一点都有\trMultiplePoint”（径改）。
\item[p.~259] $u(s,t)=u(s',t)$ 两处 $u$ 应作 $v$；又 $|s'|>\varepsilon$ $\to$ $|s'-s_0|>\varepsilon$（径改）。
\item[p.~260] $j=1,\dotsc,n$ $\to$ $j=1,\dotsc,N$（径改）。
\item[p.~260] $B_{r_1}(s_1)$ $\to$ $\overline{B}_r(s_1)$（径改）。
\item[p.~262] 命题 8.6.3 证明把 $(d\scr{F})_w(Y)$ 的一阶项统一写作 $\frac{\partial Y}{\partial s}+J_0\frac{\partial Y}{\partial t}$ $\to$ 该项本是 $\frac{\partial Y}{\partial s}+J(t,w)\frac{\partial Y}{\partial t}$ 随基点 $w$ 而变，化归 $J_0$ 须同步作\trFrame{}变换，不能保留同一 $Y$ 仅平均零阶项（结论可另补论证成立）（存疑照排）。
\item[p.~262] 归纳等式 $v(s,t)=v(k(s_\star'-s_\star),t)$ 右端 $\to$ $v(s+k(s_\star'-s_\star),t)$（照印文取 $k=0$ 即无据；取 $k\to\pm\infty$ 结论不变）（径改）。
\item[p.~264] $\Sigma=v(B_\varepsilon)$ $\to$ $\Sigma=v_2(B_\varepsilon)$（径改）。
\item[p.~270] 命题 8.7.3 前的改进不等式首项 $\|Y\|_{L^{p}(\RR\times S^{1})}$ $\to$ 应作 $\|LY\|_{L^{p}(\RR\times S^{1})}$（照印文则 $W^{1,p}$ 范数仅由两个 $L^{p}$ 范数控制，不可能成立；下文半\trFredholm{}估计的套用与 p.~284 的正确式均指此读法）（存疑照排）。
\item[p.~272] 齐次方程整体存在性称“只有当 $A$ \trMonotone{}时才有解” $\to$ \trMonotone{}性系定理 8.7.9 的充分条件而非必要（$A=-\mathrm{Id}$ 不\trMonotone{}而 $Y(s)=e^{s}Y_0$ 在全 $[0,+\infty[$ 有解）（存疑照排）。
\item[p.~274] 定理 8.7.9 “dense open subset $D(A)$” $\to$ 算子定义域通常仅要求\trDense{}（本节 $D(A)=W^{1,2}(S^1)$ \trDense{}而非开集），open 疑为衍文（径改）。
\item[p.~275] 绝对\trConvergence{}估计末行 $e^{-2\delta|\sigma|}$ $\to$ 开平方后指数应回到 $e^{-\delta|\sigma|}$（平方步的 $2\delta$ 属正当，末行疑漏开平方；\trConvergence{}结论不受影响）（存疑照排）。
\item[p.~277] $Q\cdot(AY)$ 第二项计算中 $e^{A^{-}(-\sigma)}$ $\to$ $e^{A^{-}\sigma}$（换元 $\sigma-s\mapsto\sigma$ 后不应带负号；照印文则指数增长、非\trBounded{}\trSemigroup，与 $A^{-}$ 谱为负及 $\|K(s)\|\le e^{-\delta|s|}$ 相抵）（径改）。
\item[p.~284] $W^{1,p}(\RR\times S^1,W)$ $\to$ $W^{1,p}(\RR\times S^1;\RR^{2n})$（径改）。
\item[p.~284] $s\to\infty$ $\to$ 疑应 $s\to\pm\infty$（存疑照排）。
\item[p.~288/289] $e^{tJS_{k\pm}}$、$S_\lambda^{\pm}=-J\dot\psi_\lambda^{\pm}(\psi_\lambda^{\pm})^{-1}$ 裸 $J$ $\to$ $J_0$（径改）。
\item[p.~292] 情形 $(4)$ 两维数公式的条件 $k^+>k^-$／$k^->k^+$ 与取值符号相抵 $\to$ 两条条件互换位置（径改）。
\item[p.~295] “satisfies $\widetilde{F}=0$”疑脱变元 $Z$；且 $\widetilde{F}$ 的定义域／目标域脱“；$\RR^2$”（径改）。
\item[p.~295] “the formula dor dim Ker F” $\to$ for（径改）。
\item[p.~295] $A_i(s)=\int_0^{s}-a_i(\sigma)\,d\sigma$ $\to$ $\int_0^{s}a_i(\sigma)\,d\sigma$（按 p.~293 的 $b'=-ab$、$b=e^{-A}$ 与本页 ODE $\partial Y_i/\partial s=-a_iY_i$ 及渐近式 $A_i=C+a_i^{\mp}s$，应 $A_i'=+a_i$；照原定义则 $e^{-A_i}$ 解 $Y'=+a_iY$，处处相抵）（径改）。
\item[p.~297] $f''(s)$ 第二项 $\int_0^1\bigl\langle Y,\frac{\partial S}{\partial s}\bigr\rangle dt$ $\to$ 疑漏因子 $Y$（径改）；同式首项直算应含系数 2，照排存疑（径改＋存疑并案）。
\item[p.~298] $\frac{\partial Y}{\partial s}=-J_0\frac{\partial Y}{\partial t}-SY=AY$ 之末一等号 $\to$ $-AY$（径改）；其下 $f''(s)$ 链首行 $=$ 疑应 $\ge$（径改＋存疑并案）。
\item[p.~301] $\bigl\langle Y\frac{\partial S}{\partial s},Y\bigr\rangle$ $\to$ $Y$ 与 $\frac{\partial S}{\partial s}$ 误倒，应作 $\bigl\langle S\frac{\partial Y}{\partial s},Y\bigr\rangle$（径改）。
\item[p.~301] $\frac{1}{r^2}\int_{B_r(s_0,t_0)}w$ $\to$ $\frac{1}{\pi r^2}\int_{B_r(s_0,t_0)}w$（径改）。
\item[p.~302] 分母 $(1-\rho)^2$ $\to$ $(1-\rho_\star)^2$（径改）。

\subsection*{第~9~章}
\item[p.~307] $(x,\sigma)\mapsto u(s+\sigma,t)$ $\to$ 首分量应作\trTrajectory{} $u$，且与端点 $x$ 相区别（径改）。
\item[p.~308] 并集下标 $\mu(y)=\mu(x)+1$ $\to$ $\mu(y)=\mu(x)-1$（同页 $\sum$ 下标同病，一并径改）（径改）。
\item[p.~309] 定理 9.1.7 自 $u^{0}\in\scr{M}(x_{0},x_{1})$、$(s_{n}^{k})_{0\le k\le\ell}$ 起（$0$ 起标），而证明先构 $u^{1}$、归纳又写 $(s_{n}^{0}),\dots,(s_{n}^{k})$ 且 $(s_{n}^{0})$ 于证明中从未定义 $\to$ 两套起标错开一位，应统一（存疑照排）。
\item[p.~310] $u_n\cdot s_n^{k+1}\in B(x_k,\varepsilon)$ 未带求值记号 $\to$ 疑应补出 $(s)$（存疑照排）。
\item[p.~311] $\partial\overline{\scr{L}}(x,y)$ $\to$ $\partial\overline{\scr{L}}(x,z)$；右端下标 $\mu(x)<\mu(y)<\mu(z)$ $\to$ $\mu(z)<\mu(y)<\mu(x)$（径改）。
\item[p.~315] $\exp_{y(t)}^{-1}\bigl(y(s-\rho,t)\bigr)$ $\to$ $v(s-\rho,t)$（径改）。
\item[p.~315] 性质 (3) 仅以 $Y$、$Z$ 为 $\cC^{\infty}$ 类为据断言 $\cC^{\infty}_{\mathrm{loc}}$ \trConvergence $\to$ \trSmooth{}性并不提供诸阶导数的衰减或 $L^{p}$ 控制（尖峰平移反例），须补各阶导数的相应条件（存疑照排）。
\item[p.~319] “当 $\rho$ 趋于无穷时 $S_{\rho}$ 趋于……$x$ 与 $z$” $\to$ 两端极限实为 $s\to\pm\infty$（径改）。
\item[p.~320] 引理 9.4.4 仅设 $G$ 连续 $\to$ 证明反复使用 $G$ 的可加性（仅有连续性时有显式反例），须设连续线性（径改）。
\item[p.~324] $\sup$ 下标 $z\in[x,y])$ 括号多一枚 $\to$ 径删（径改）。
\item[p.~325] $(dF)_{w_{\rho}}(Y)$ 中 $(dJ)$ 项印作负号 $\to$ 对 $w_{\rho}+\varepsilon Y$ 直接求导应为正号（径改）。
\item[p.~327/329] $\beta^{-}(s+1)$ 的过渡区实为 $[-2,-1-\varepsilon]$：原书用作控制区间的 $[-1,1]$ 与\trSupport{} $[-1,0]$ 均应扩至 $[-2,0]$ 一类区间（存疑照排）。
\item[p.~331/332] 命题 9.4.15 证明中 $\bigl\|G_{\rho}(d\scr{F}_{\rho})_{\gamma(\rho)}-\mathrm{Id}\bigr\|$ 的小性只在子空间 $W_{\rho}^{\perp}$ 上成立，却被作用于环境向量 $\frac{\partial\gamma}{\partial\rho}$（$\gamma\in W_{\rho}^{\perp}$ 不蕴涵其导数属之：\trComplement{}随 $\rho$ 移动）$\to$ 照排，注载修补路径：取沿 $W_{\rho}^{\perp}$ 到 $W_{\rho}$ 的 $L^{2}$ 投影 $P_{\rho}$，对 $P_{\rho}\gamma=0$ 求导后对 $Q_{\rho}\frac{\partial\gamma}{\partial\rho}$ 用原逆估计，$P_{\rho}'$ 的一致界由两平移核向量的 Gram 矩阵趋于\trNonsingular{}对角而得，结论不变（存疑照排）。
\item[p.~332] Conclusion 段“$\psi(\rho)$ 与它的导数都在 $W^{1,p}$ 中趋于 $0$” $\to$ $\psi(\rho)$ 是取值 $\scr{M}(x,z)$ 的解、平移后\trConvergence{}于 $u$ 与 $v$，趋 $0$ 者是 $\gamma$ 与 $\gamma'$（径改）。
\item[p.~335] $w_{\nu,\rho}$ 末行 $v(s-\nu(\rho))$ $\to$ $v(s-\nu(\rho),t)$（径改）。
\item[p.~337] $\exp_{w_{\nu,\rho_{n}}}\gamma_{n}(h)$ 之 $(h)$ 并非衍字：$\gamma_{n}(h)$ 系“薄片 $h+W_{\rho_{n}}^{\perp}$ 上的\trNewton--\trPicard{}解”这一参数族的记法，前文无自变元的 $\gamma_{n}$ 即零薄片解 $\gamma_{n}(0)$（照排备案；外审二轮改判）。
\item[p.~341] $\scr{A}_H(\lambda_n(s_n+\sigma,\cdot)\le\scr{A}_H(\tau,\cdot))$ $\to$ 左端补右括号、右端补出 $\lambda_n$（径改）。
\item[p.~341] $\lambda=\lim\ell_n$ $\to$ $\lambda=\lim\widetilde{\ell}_n$（径改）；同段 $\lambda\notin\{\widehat{u},\widehat{v}\}$ 空间不配，照排存疑（径改＋存疑并案）。
\item[p.~344] $\scr{A}_H\bigl(\widetilde{\ell}_n(\sigma'_n)\bigr)$ $\to$ $\scr{A}_H\bigl(\widetilde{\ell}_n(-\sigma'_n)\bigr)$（径改）。
\item[p.~344] $[-\sigma'_n,\sigma_n]$ $\to$ $[-\sigma'_n,-\sigma_n]$（末端 $\sigma_n$ 疑漏一枚负号；径改）。
\item[p.~344] 引理 9.6.12 证明\trEnergy{}差等式第二行重积分 $\int_{-\sigma_n}^{-\sigma'_n}$ 积分限反向 $\to$ 互换作 $\int_{-\sigma'_n}^{-\sigma_n}$（左端系作用递减所得的正\trEnergy{}差，以反向积分限积非负被积函数其值非正，二者相抵；径改）。
\item[p.~345] 下界 $\gamma_0^2(\sigma'_n-\sigma_n)$ $\to$ $\gamma_0(\sigma'_n-\sigma_n)$（径改）。
\item[p.~355] $d(dF)_{\gamma_n(h)}$ $\to$ $(dF)_{\gamma_n(h)}$（径改）。
\item[p.~349/351] 两处 $\int e^{\delta s}\,ds$ 等号右端均漏因子 $1/\delta$（后续估计只消把 $1/\delta$ 并入常数即恢复原状）（存疑照排）。
\item[p.~350] 以\trFloerEquation{}替换处两括号内印作 $X_{t}(w_{n})-J(w_{n})\cdot\frac{\partial w_{n}}{\partial s}$ 一类 $\to$ 按 $\scr{F}(w)=0$ 应解出 $X_{t}(w)+J(w)\cdot\frac{\partial w}{\partial s}$，均应为正号（存疑照排）。
\item[p.~351] $\varGamma$ 链首行第二项 $\varGamma(w_n(s,t),\widetilde{\ell}_n(s,t))$ $\to$ $\varGamma(w_n(s_1,t),\widetilde{\ell}_n(s_1,t))$（径改）。
\item[p.~353] 引理 9.6.16 末 $\|\gamma_{h}\|\le\|GF(0)\|+\|h\|$ $\to$ \trContractingMap{}估计只能给 $2(\|GF(0)\|+\|h\|)$（数值反例可直见），且 $G$ 须设连续线性（存疑照排）。
\item[p.~356] 引理 9.6.18 证明 $\widehat{\lambda}_{k}\in[a-\delta_{1},a-\delta_{2}]$ $\to$ $I_{\varepsilon_{k}}\supset[a,1[$ 不保证越过事先取定的 $a-\delta_{2}$，最小修补径取 $\widehat{\lambda}_{k}=a$（$a<c$；其后两段反证照旧）（存疑照排）。

\subsection*{第~10~章}
\item[p.~363] \trLinearization{}矩阵的分量式 $\sum_{j}(a_{ji}+S_{ji})\,y_{j}$ 与“令 $A_{ij}=a_{ij}+S_{ij}$”（列向量 $AY$）不一致 $\to$ 应作 $A_{ij}=a_{ji}+S_{ji}$（或把\trFrame{}系数定义与分量式一并转置；并宜先作切向投影再取切向系数）（存疑照排）。
\item[p.~365] $\|Y\|_{W^{1,2}}=\int\cdots ds$ 左端漏一枚平方 $\to$ $\|Y\|_{W^{1,2}}^2$（径改）。
\item[p.~367] $\frac{1}{C_1}\|Y\|_{W^{1,2}}$ 同漏平方 $\to$ $\frac{1}{C_1}\|Y\|_{W^{1,2}}^2$（径改）。
\item[p.~367] 引理 10.2.4 证明 $|\langle(iv\cdot\mathrm{Id}+B)V,V\rangle|\ge\sqrt{v^{2}+C_{0}^{2}}\,\|V\|^{2}$ 对不定对称 $B$ 不真（反例 $\operatorname{diag}(1,-1)$）$\to$ 结论由\trOrthogonal{}特征基直接计算 $\|(iv\cdot\mathrm{Id}+B)V\|^{2}=\sum_{j}(v^{2}+\lambda_{j}^{2})|V_{j}|^{2}$ 挽救（存疑照排）。
\item[p.~368] “The cokernel of $u$” $\to$ “The cokernel of $L_u$”（径改）。
\item[p.~369] “when $L_u$ connects two nondegenerate critical points” $\to$ $u$（径改）。
\item[p.~370] $Y_0(s)/ds=-AY_0(s)$ $\to$ $\frac{dY_0}{ds}(s)=-AY_0(s)$（径改）。
\item[p.~373] $(d\scr{G})_{(u,J_0)}$ 下标及定义域第二因子 $\scr{J}_0(\delta)$ $\to$ 疑应 $(u,J)$、$T_J\scr{J}_0(\delta)$（存疑照排）。
\item[p.~374] $\varGamma$ 的像印作“有限维的闭子空间” $\to$ $\operatorname{Im}\varGamma\supseteq\operatorname{Im}L_{u}$ 而后者闭且余维数有限，应作“余维数为有限的闭子空间”（径改）。
\item[p.~375] 引理 10.3.5 证明的块形 $[\,[A,B],[-B,A]\,]$（$B$ \trAntisymmetric）与反 $\CC$-线性相抵 $\to$ 应作 $[\,[A,B],[B,-A]\,]$（$A$、$B$ 对称）（存疑照排）。
\item[p.~377] 引理 10.4.1 证明的\trOrientation{}积分限 $\int_{t}^{t_{1}}$ 当 $t_{1}<t$ 时为反向积分、相应估计失效 $\to$ 应以 $\min/\max$ 为限并配合\trHoelder{}不等式（存疑照排）。
\item[p.~380] \trBrokenTrajectory{} $(\widehat{u},\widehat{v})$ $\to$ 按上下文疑应 $(\widehat{v},\widehat{w})$（存疑照排）。
\item[p.~381] 证明末句 “contradicts the fact that $v_{n_{k}}$ does not depend on $t$” 多一枚 not $\to$ 按归谬设定，矛盾应针对“$v_{n_{k}}$ 依赖于 $t$”（径改）。

\subsection*{第~11~章}
\item[p.~386] $X_{t,s}$／$dH_{t,s}$ $\to$ $X_{s,t}$／$dH_{s,t}$（径改）。
\item[p.~387] 误差估计取 $k=\sup\frac{\partial H}{\partial s}$ 照排——$k$ 漏了积分区间 $[-R,R]$ 的长度因子，应作 $k=2R\max\bigl(0,\sup\frac{\partial H}{\partial s}\bigr)$（取 $H_{s,t}=H_{0}+\epsilon s$ 即成反例）（存疑照排）。
\item[p.~388] “此邻域内 $\scr{H}_{\mathrm{reg}}$ 的诸稠开子集的一个可数交”与第 8 章定理 8.1.1 的口径不合，照排存疑；“此外”句 $u\in\scr{M}^\varGamma(x,y)$／$(d\scr{F}^\varGamma)_u$ 疑应带上标 $(h)$，径改（存疑＋径改并案）。
\item[p.~388/415] $\cC^{\infty}_{\varepsilon}$（紧\trSupport{}＋$\|h\|_{\varepsilon}<+\infty$）在范数 $\|\cdot\|_{\varepsilon}$ 下不完备（取紧支的 $b$，其诸平移 $b(\cdot-j)$ 加权和即给出极限\trSupport{}无界的\trCauchy{}列），而本章\trTransversality{}论证需要完备性；§11.3.b（p.~415）的同类定义同病 $\to$ 两处照排，注载修补路径（固定共同紧\trSupport{}或改完备的衰减空间，均须重写端点论证）（存疑照排）。
\item[p.~389] $W^{1,p}(u^\star W)$／$L^p(u^\star W)$ $\to$ $u^\star TW$（径改）。
\item[p.~390/395/409/411] 四小项归组：$\operatorname{Supp}(h)\cap\operatorname{Im}(\widetilde{u})=\widetilde{u}(C_{\delta})$ 之等号无据（只能得 $\subseteq$）；破碎链两条各多写一项（末端用到无定义的 $y_{\ell+1}$、$x_{k+1}$，应作 $j\le\ell-1$、$i\le k-1$）；$[\sigma_{n},\sigma'_{n}]\subset[A,B]$ 系反序空区间 $\to$ 径改作 $[\sigma'_{n},\sigma_{n}]$；$\gamma_{n}$ 的自变元印作 $h_{n}(s)$ $\to$ 应接收整\trSection{} $h_{n}$（存疑照排；其三径改）。
\item[p.~405] $\scr{A}_{H^b}\bigl(\lambda_n(\sigma_n+\tau)\bigr)$ $\to$ $\lambda_n(\sigma_n+\sigma)$（径改）。
\item[p.~406] 引理 11.2.6 未排除破碎极限尾端为常值 continuation 分量的情形——此时等临界作用量时刻不存在、证明中的三球选择无从进行（原书 p.~444 注记 11.5.7 自认常值解不可排除）（存疑照排）。
\item[p.~411] $h+W_u^\perp\cap B(0,\varepsilon_1)$ $\to$ $W_{\rho_n}^\perp$（径改）。
\item[p.~412] $\gamma_n$ 自变元 $h_n$ $\to$ $h$（径改）。
\item[p.~412] $I_\varepsilon$ 并集各项 $\exp_{\rho_n}(\cdots)$ $\to$ $\exp_{w_{\rho_n}}(\cdots)$，且各项漏写平移，一并补全为 $(s,t)\mapsto\exp_{w_{\rho_n}(s-\rho_n,t)}\bigl(i_{\rho_n}(\gamma_{n}(h))(s-\rho_n,t)\bigr)$（径改）。
\item[p.~414/419/423] 下标 $\scr{A}_{H_b}$（本节九处） $\to$ $\scr{A}_{H^b}$（径改）。
\item[p.~417] “为证定理 11.3.3” $\to$ 按上下文应为定理 11.3.2（径改）。
\item[p.~417] \trTangentSpace{}方程两处漏 $-(\partial J^{\lambda}/\partial\lambda)X$ 项（同病另见 p.~428 一处、p.~437 一处，彼处更把标量 $\partial H/\partial\lambda$ 直接加于向量，应为其\trGradient）$\to$ 按 $\operatorname{grad}^{J}H=-JX_{H}$ 对 $\lambda$ 求导补齐，四处一并径改（径改）。
\item[p.~417/418] 命题 11.3.5 称 $\pi$ 的\trRegularValue{}之集是“$0$ 的一个邻域内的\trDense\trOpen{}子集”照排——\trSard--\trSmale{}只给残余性，“$0$ 的邻域”与\trOpen{}性均超出（\trOpen{}性须另证且依赖诸解的\trEnergy{}一致界），下文“由这条命题所给出的那个邻域”沿用同说（同第 8 章 §8.5 末口径）（存疑照排）。
\item[p.~422] 首式求和范围印作 $y\in\operatorname{Crit}\scr{A}_{H^{a}}$ $\to$ $m^{\varGamma}(x,y)$ 所数为 $H^{a}\to H^{b}$ 的延拓解，$y$ 应属 $\operatorname{Crit}\scr{A}_{H^{b}}$（下一页展开同一求和即已印 $H^{b}$）（径改）。
\item[p.~420] $(u,\lambda)\in\scr{M}^{\varGamma_\lambda}$——按定义其元素是解 $u$ 而非偶对（存疑照排）。
\item[p.~420] 常数 $C$ 的定义式第二下标 $\operatorname{Crit}H^{b}$ 与同式第一下标 $\operatorname{Crit}\mathscr{A}_{H^{a}}$ 体例不对称，疑漏排 $\mathscr{A}$（字面尚可两读）（存疑照排）。
\item[p.~421] “tends to infinity” $\to$ 按上下文应为一 $\infty$（径改）。
\item[p.~424] $\scr{M}^\varGamma(x,z)\cup\varPi^\varGamma(x,z)$ 的紧性句首因子 $\scr{M}^\varGamma(x,y)$ $\to$ $\scr{M}^\varGamma(x,z)$（径改）。
\item[p.~426] $\lim_{s\to+\infty}\psi(s,\bullet)=z$ 之 $\psi$ 无 $\rho$ 下标 $\to$ $\psi_\rho$（径改）。
\item[p.~428] $L^\varGamma_\rho(a,0)$ 计算第二行 $J\cdots\frac{\partial w}{\partial t}$ $\to$ $\frac{\partial w_\rho}{\partial t}$（径改）。
\item[p.~429] \trIndex{}路径定义式左端 $L_{\rho,\sigma}$ $\to$ $L_{\rho,\sigma}^\varGamma$（径改）。
\item[p.~429] $\scr{Y}(a,Y)=aV+L^u(Y)$ 之 $\scr{Y}$ 无上标 $\to$ $\scr{Y}^u$（径改）。
\item[p.~431] $\dot\beta^-Y_n(s+1,t)$ $\to$ $\dot\beta^-(s+1)Y_n(s,t)$（径改）。
\item[p.~432] 两处极限下标 $\rho\to0$ $\to$ $\rho\to+\infty$（径改）。
\item[p.~436] “for every $u$, a vector $Y_n$” $\to$ “for every $n$”（径改）。
\item[p.~438] 命题 11.5.2 印作 $\Phi^{\varGamma''}\circ\Phi^{\varGamma''}$ $\to$ 按两因子各自的定义域与陪域应作 $\Phi^{\varGamma''}\circ\Phi^{\varGamma'}$（径改）。
\item[p.~440] “the injectivity of $\varphi$” $\to$ $\chi$（径改）。
\item[p.~445] $u=\lim$ 式首项下标 $s\to+\infty$ $\to$ $n\to+\infty$（径改）。
\item[p.~445] “不会有多于两条在 $s$ 方向上为常值的\trTrajectory” $\to$ 疑应“非常值”（存疑照排）。
\item[p.~448] $\ell(0)\in B(y,\delta)$ $\to$ $\ell_n(0)$（径改）。
\item[p.~448] $\ell_n(\sigma'_n)\in\partial B(x,\delta)$ $\to$ $\ell_n(-\sigma'_n)$（径改）。
\item[p.~450] $Y_{n}\in W^{1,p}(w_{\rho}^{\star}TW)$（下标单 $\rho$）$\to$ 与命题陈述处的 $w_{\rho_{n}}^{\star}TW$ 相对应，下标疑脱一 $n$（径改）。
\item[p.~451] 同一复合式又印作 $\Phi^{\varGamma_{\rho}}=\Phi^{\varGamma'}\circ\Phi^{\varGamma''}$ $\to$ 右端应作 $\Phi^{\varGamma''}\circ\Phi^{\varGamma'}$（与 p.~438 处的径改相一致）（径改）。

\subsection*{第~12~章}
\item[p.~454] $C(k,p,K,U)$ $\to$ $C(k,p,V,U)$（径改）。
\item[p.~455] “the equality holds” $\to$ “the inequality”（径改）。
\item[p.~455] 引理 8.7.2 转述的加强断言“更一般地，若 $Y$、$LY\in L^{p}$ 则 $Y\in W^{1,p}$”未获证明（本节证明恒以 $Y\in W^{1,p}$ 为前提）$\to$ 须先局部\trRegular{}化、平台截断用统一估计再令截断趋于 $\infty$ 补足（存疑照排）。
\item[p.~456] $\varphi$ “is continuous” $\to$ 其后用到 $(d\varphi)_u$，疑应 $\cC^1$ 类（径改）。
\item[p.~470] 定理 12.1.2 证明末 $u\in W^{1,p/2}(\RR\times S^{1})$ $\to$ 与同页同款 $W^{2,p}$ 处一致，疑漏 $\mathrm{loc}$（径改）。
\item[p.~458] 定理 12.1.5 陈述两处 $W^{\min(k,1),p}$（$\widetilde{J}$ 假设及其界 $c_{0}$）$\to$ 按 §12.4.c 的证明应为 $W^{\max(k,1),p}$；\trIndex{}集 $\ell\in\{0,\dots,\max(k-1,0)\}$ $\to$ 按该证明（对每个 $\ell\le k$ 作估计）与结论 $u\in W^{k+1,p}$ 应作 $\{0,\dots,k\}$（两项并案径改）。
\item[p.~459] $\|Y\|$ 项范数下标 $L^p_{\times S^1}$（两处） $\to$ $L^p(\RR\times S^1)$（径改）。
\item[p.~459] “therefore the desired equality” $\to$ “the desired inequality”（径改）。
\item[p.~460] \trPlateauFunction{}条件 “zero for $|s|\le M+1$ and is $1$ for $|s|\le M$” $\to$ 第一个不等号应作 $|s|\ge M+1$（径改）。
\item[p.~462] 命题 12.3.1 证明两处估计缺口并案：\trConvolution{} $(\partial\psi/\partial s)\star u$ 在 $V$ 上取值越出 $V$，仅以 $\|u\|_{L^{p}(V)}$ 控制不足 $\to$ 应取稍大集（含于 $U$）上的 $L^{p}$ 范数；由 $\|\nabla v\|_{L^{p}}$ 过渡 $\|v\|_{W^{1,p}}$ 所用\trPoincare{}不等式需零均值条件而 $v$ 未必满足 $\to$ 可再经 Young 不等式逐项控制修复（存疑照排）。
\item[p.~463] $\frac{\partial(\beta g)}{\partial t}$ $\to$ $\frac{\partial(\beta h)}{\partial t}$（径改）。
\item[p.~464] 归纳步叙述印作“假设对 $k$ 成立、证对 $k+1$” $\to$ 下文实设 $u,f\in W^{k,p}_{\mathrm{loc}}$ 行 $k$ 层结论、所用归纳假设系 $k-1$ 层，\trIndex{}错位一格（径改）。
\item[p.~467] “the equality of Theorem 12.1.5” $\to$ “the inequality”（径改）。
\item[p.~468] 第一步 $\Delta u$ 展开式第二、三项印作 $\cdot u$ $\to$ 按乘积法则应为 $\cdot\frac{\partial u}{\partial t}$（引理 12.4.3 证明中的同款展开即作 $\frac{\partial u}{\partial t}$，无误）（存疑照排）。
\item[p.~470] 估计式中 $(1+kc_{0})$ 的 $k$ 在该证明中无定义 $\to$ 应即\trSobolev\trEmbedding{}中的常数 $K$（或径吸入新常数 $c$）（存疑照排）。
\item[p.~471] 性质条 $q_{0}=p$ 与引理 12.4.3 前提 $1/p+1/q_{0}=1/r_{0}$ 不\trCompatible（强制 $r_{0}=p/2$，与上文序列 $r_{0}\in\,]1,2[$、$r_{0}\le p/2$ 相抵；$p\ge 4$ 时又不在 $]1,2[$ 内）$\to$ 注载修补：取 $r_{0}\in\,]1,\min(2,p/2)]$、按关系定 $q_{0}\le p$ 并用局部 $L^{p}\subset L^{q_{0}}$（存疑照排）。
\item[p.~472] $\|u\|_{W^{1,p}(U)}$ $\to$ $\|u\|_{W^{1,p}(V)}$（径改）。
\item[p.~473] $u\in W^{2,p}(\RR\times S^1)$ $\to$ 补 $\mathrm{loc}$（径改）。
\item[p.~473] $\widetilde{h}\in W^{1,k+1}_{\mathrm{loc}}$ $\to$ 波纹号衍（应作 $h$）、上标应作 $k+1,p$（径改）。
\item[p.~473] $r_0=p/2$ $\to$ $r_0\le p/2$（径改）。
\item[p.~474] 引理 12.4.5 证明及下文对仅相对紧的 $V$、$U$ 直接用 $W^{1,p}\subset L^{\infty}$ \trEmbedding（通常需边界\trRegular{}或可延拓域）$\to$ 注载补救：在 $\overline{V}$ 外取有限个\trSmooth{}小域作内估计（存疑照排）。
\item[p.~475] 本证明之 $\ell\in\{0,\dots,k\}$ 本身不误，与 §12.1 陈述处原书所印 $\max(k-1,0)$ 相抵——原书之误在陈述处，已径改（见 p.~458 条）。

\subsection*{第~13~章}
\item[p.~481] $i_{\rho}$ 之 $\cC^{1}$ 一致界的原书理由（参数集 $[\rho_{0},+\infty[\times[-1,1]\times S^{1}$ “\trCompact{}”）不成立 $\to$ 改由 §9.4 的 $\cC^{\infty}_{\mathrm{loc}}$ \trConvergence{}（充分大 $\rho$）＋\trBounded{} $\rho$ 区间 $[\rho_{0},R]$ 的\trCompact{}性得界（注中并明 $\cC^{1}$ 只含 $s,t$ 导数）（径改）。
\item[p.~483/484] $P$ 定义链三式：$\overline{g}$ 取于起点 $p$（且 $\widetilde{\scr{T}}$ 脱基点）、后两槽次序与分量展开相反 $\to$ 统一改在 $q(p,X)=\exp_{p}(\pi_{p}X)$ 处取\trMetric、第三槽\trFrame{}／第四槽残差，零式作 $P(p,0,A,B)=\overline{g}_{p}(\pi_{p}A,B)$（径改）。
\item[p.~485] $\bigl\|(d\psi)_Z\circ\varphi-(d\psi)_0\circ\varphi\bigr\|^{\mathrm{op}}$ $\to$ $\bigl\|d(\psi\circ\varphi)_Z-d(\psi\circ\varphi)_0\bigr\|^{\mathrm{op}}$（径改）。
\item[p.~485] 链式第二行第二项 $(d\varphi)_0(v)$ $\to$ $(d\psi)_{\varphi(0)}\bigl((d\varphi)_0(v)\bigr)$（径改）。
\item[p.~484] “We suppose there that exist…”之 there that 语序滑写 $\to$ 显见排印笔误，径顺不改义（径改）。
\item[p.~489] $\psi=\mathrm{Id}\,\widetilde{\scr{F}}_\rho$ $\to$ $\psi=\mathrm{Id}\times\widetilde{\scr{F}}_\rho$（径改）。
\item[p.~489] 靶空间 $L^p(\RR\times S^1;\RR^m)$ $\to$ $\RR^{2n}$；分量范围 $i=1,\dots,n$ $\to$ $i=1,\dots,2n$（径改）。
\item[p.~489] $X=0$ 四项式把 $D_{2}\widetilde{\exp}(w_{\rho},0)v$ 直接写成任意环境向量 $v$（实为 $\pi_{w_{\rho}}v$）$\to$ 按 $V=\varPi_{\rho}v$ 改写并补 $D\pi$ 零阶项（所需 $W^{1,p}$ 算子范数估计不变）（径改）。
\item[p.~490] 引理 13.4.2 左端第二项印作 $(dp_{\rho})_{0}$ $\to$ 右端所量为 $(X,Y)$ 到 $\varphi(0)$ 的距离（照字面取 $(X,Y)=\varphi(0)$ 则右端为 $0$ 而左端一般非 $0$），应作 $(dp_{\rho})_{\varphi(0)}$（径改）。
\item[p.~492] $\bigl\|\frac{\partial\varphi_{\rho}}{\partial\rho}(0)\bigr\|\le k_7$ 之分子原书印作 $\varphi_{\rho}$（条件 (3) 第二项分子同病印作 $\psi_{\rho}$）$\to$ 疑漏偏导号 $\partial$，径改作 $\frac{\partial\varphi_{\rho}}{\partial\rho}$（径改）。
\item[p.~492] 估计链末行括号内末项印作 $k_{7}$ $\to$ 按上一行 $k_{3}\|\varphi_{\rho}(z)-\varphi_{\rho}(0)\|(\cdots)$ 以假设 (2)、(4) 代入展开应作 $k_{3}k_{7}$（以统一常数 $M$ 重新吸收后结论不变）（存疑照排）。
\item[p.~494] $\frac{\partial\widetilde{\scr{F}}_{\rho}}{\partial\rho}$ 展开式的两个 $D_{1}\widetilde{J}$ 项首变元均直接印作 $\frac{\partial w_{\rho}}{\partial\rho}$ $\to$ 按链式法则应经 $D_{1}\widetilde{\exp}_{(w_{\rho},X)}\bigl(\frac{\partial w_{\rho}}{\partial\rho}\bigr)$（末项 $(d\operatorname{grad}H_{t})$ 中恰保留该因子，可资内证），补入后估计照常成立，两处径改（径改）。
\item[p.~497] $D_2\widetilde{\exp}_{(w_\rho,\gamma(\rho))}\bigl(\frac{\partial w_\rho}{\partial\rho}\bigr)$ $\to$ 变元应作 $\bigl(\frac{\partial\gamma}{\partial\rho}\bigr)$（径改）。
\item[p.~498] 极限号下 $s\to+\infty$ $\to$ $\rho\to+\infty$（径改）。
\item[p.~498] “对 $s\ge-1$” $\to$ 按 $w_\rho$ 的分段定义应作 $s\le-1$（径改）。
\item[p.~502] 积分区域 $]-\infty,1]\times S^1$ $\to$ $]-\infty,-1]\times S^1$（径改）。
\item[p.~504] 复述 $p$ 的定义时标量映射 $P$ 的陪域印作 $\RR^{2n}$ $\to$ $P$ 系逐分量使用的实值映射（彼处陪域即 $\RR$），应作 $\RR$（径改）。
\item[p.~505] 引理 13.7.1 的型别印作 $W^{1,p}(\RR\times S^{1};\RR^{2n})$ $\to$ 按上文 $\scr{N}\colon W^{1,p}(\RR\times S^{1};\RR^{m})\to L^{p}(\RR\times S^{1};\RR^{m})$ 应作 $\RR^{m}$（径改）。
\item[p.~508/509/512] 参数方向场三处漏 $-(\partial J/\partial\lambda)X_{H}$ 型项——$d\overline{F}_{\rho}$ 的 $\lambda$ 导数展开式（原书 p.~508）、$\widetilde{V}_{\rho}(\lambda+\lambda_{\star},Z)$ 定义式及其展开（原书 p.~509）与 $V_{\rho}(\lambda+\lambda_{\star},0)$ 定义式（原书 p.~512），与第 11 章同一疏漏同源 $\to$ 三处一并径改（径改）。
\item[p.~508] 方向导数取值点印作 $\sigma=\lambda$ $\to$ 被导函数为 $\sigma\mapsto\scr{F}_{\rho}^{\varGamma_{\lambda+\lambda_{\star}+a\sigma}}$，下一行所得恰是 $\sigma=0$ 处的链式结果，应作 $\sigma=0$（下一行诸式不受影响）（径改）。
\item[p.~511] $\widetilde{V}_\rho(s,t)=\cdots$ $\to$ $\widetilde{V}_\rho(0)(s,t)$，$\widetilde{\exp}$ 下标随之为 $(w_\rho,0)$（径改）。
\item[p.~511] 引理 13.8.2 证明只控 $\bigl\|(d\chi)_{ra+(1-r)b}-(d\chi)_{b}\bigr\|^{\mathrm{op}}$，而均值不等式所需是 $\sup_{r\in[0,1]}\bigl\|(d\chi)_{ra+(1-r)b}\bigr\|^{\mathrm{op}}$ 本身——缺基线一项（第二行等号处又消失一枚因子 $r$）$\to$ 照排，注载修补：另估 $\bigl\|(d\chi)_{b}\bigr\|^{\mathrm{op}}\le k_{7}$ 后相加，最终结论仍成立（存疑照排）。
\item[p.~513] $\bigl(V_\rho(\lambda+\lambda_\star),0-V_\rho(\lambda_\star,0)\bigr)$ $\to$ $V_\rho(\lambda+\lambda_\star,0)-V_\rho(\lambda_\star,0)$（径改）。

\subsection*{第~14~章}
\item[p.~521/522] 练习 18--21 与练习 34 的归一化并案照排：按 $\exp(2\pi i t)$\trParametrization{}的圆作用要求时刻 $1$ 流为恒等，而所给\trHamiltonian{}量的流为 $t\mapsto[e^{im_{i}t}z_{i}]$（差 $2\pi$ 因子）；同一约定下 $\CP^{1}(\CC)$ 的面积大小为 $\pi$，$\int_{S^{2}}w^{\star}\omega$ 只能落在 $\pi\ZZ$ 而非 $\ZZ$（存疑照排）。
\item[p.~522/523] 练习 23 与练习 27 照排——末一断言须以 $W$（相应地 $V$）\trCompact{}（或\trTrajectory{}相对紧）为前提，非紧时有反例（$\CC^{2}\setminus\{0\}$ 配 Hopf 圆作用无\trFixedPoint；$\RR$ 上 $d(\arctan x)$ 的\trTrajectory{}的\trEnergy{}有限而不连零点）（存疑照排）。
\item[p.~523] “\trHamiltonian\trVectorField{} $X_H$ 与 $H_K$”之 $H_K$ $\to$ $X_K$（径改）。
\item[p.~523] 练习 25 末段照排两点：生成 $\sigma_{t}=\psi_{t}\circ\varphi_{t}$ 的 $F_{t}=K+H\circ\psi_{t}^{-1}$ 一般非时间 $1$ \trPeriodic，须先作\trPeriodic{}延拓并区分闭区间\trTrajectory{}与整\trPeriodic{}解；$\sigma_{1}$ 的 $16$ 个\trFixedPoint{}仅 $4$ 个\trNondegenerate{}且恰为常值（\trContractible），“$F_{t}$ 是\trNondegenerate{}的”只在\trContractible{}一类上成立（存疑照排）。
\item[p.~524] 练习 28 并案：首行漏靶空间，径补 $\longrightarrow W$（径改）；$\widetilde{u}(s,t)=(\psi^{t})^{-1}(u(s,t))$ 在 $t$ 上一般非\trPeriodic（满足扭曲边界条件 $\widetilde{u}(s,t+1)=(\psi^{1})^{-1}(\widetilde{u}(s,t))$；与之\trCompatible{}的取法为 $\widetilde{J}_{t}=(\psi^{t})^{\star}J$、$\widetilde{H}_{t}=(H_{t}-K_{t})\circ\psi^{t}$，或补设 $\psi^{1}=\mathrm{Id}$）照排（存疑＋径改并案）。
\item[p.~525/526] 练习 34 照排——$\varphi_{\alpha}$ 的像一般只是 $\ZZ$ 的子群 $d\ZZ$，$\exp\colon\RR\to\RR/\ZZ$ 沿 $f$ 的拉回共有 $d$ 个\trConnectedComponent（$\alpha=0$ 时拉回即 $V\times\ZZ$），覆叠群为 $\ZZ$ 者仅是含指定\trLift{}基点的那个分支（或补设 $\varphi_{\alpha}$ 为满）（存疑照排）。
\item[p.~527] “the stabilizer the last vector” $\to$ “the stabilizer of the last vector”（径改）。
\item[p.~527] 分块矩阵左下块 $-\alpha$ $\to$ $\alpha$（径改）。
\item[p.~527/528] 练习 39 三点并案照排：显式 $4\times4$ 矩阵诸 $2\times2$ 块的旋转方向与 $\exp(t\alpha)$ 相反（相应于 $-\alpha$）；(2) 的 $X=(1,-i,0,0)$ 只与所印矩阵\trCompatible；“当 $t\neq0$ 时属 $\mathrm{Sp}(2n)^{+}$”于 $t\in2\pi\ZZ$ 失效（应限 $t\notin2\pi\ZZ$）——$\rho(A)=1$ 与方向取舍无关（存疑照排）。
\item[p.~528] \trLagrange{}式\trImmersion{} $f\colon L\to\RR^n$ $\to$ $\RR^{2n}$（径改）。
\item[p.~529] 练习 42 末段照排——$\pi_{2}(W)=0$ 只保证 $j_{\star}\colon\pi_{1}(L)\to\pi_{1}(W)$ \trInjective，$\mu_{L}$ 的定义域应限 $\operatorname{Ker}[\pi_{1}(L)\to\pi_{1}(W)]$（或补设 $\pi_{1}(W)=0$；$W=T^{\star}\mathbf{T}^{2}$、$L$ 为\trZeroSection{}即反例）（存疑照排）。
\item[p.~530] $\delta$＝$|P|$ 在 $\operatorname{Im}\gamma$ 上的上确界 $\to$ 欲使结论可证疑应取下确界（存疑照排）。
\item[p.~530] $h$ 未曾定义 $\to$ 按文义疑应 $h=Q/P$（存疑照排）。
\item[p.~530] display 第三项 $\sigma\cdot Y$ $\to$ $\Sigma(s)\cdot Y$（径改）。
\item[p.~531] 练习 46 算子定义印作 $Y\mapsto\frac{dY}{ds}+A_{0}(s)$ $\to$ 向量与矩阵不可相加，应作 $Y\mapsto\frac{dY}{ds}+A_{0}(s)Y$（补后 $\operatorname{ind}L_{0}=m_{+}-m_{-}$ 与命题 10.2.8 的 Morse \trIndex{}之差相合）（径改）。

\subsection*{附录~A}
\item[p.~539] §A.1.g 首句印作 “an exterior derivative of degree $p$” $\to$ 按上下文实指 $T_{x}V$ 上的交替 $p$-线性形式（$p$ 次外形式），非外\trDifferential{}算子 $d$（存疑照排）。
\item[p.~546] $\scr{U}'_a=\{g\in\scr{U}_a\mid f\pitchfork P$ 沿 $K_a\}$ 及下一段行文之 $f$ $\to$ 疑应指 $g$（存疑照排）。
\item[p.~546] “作为稠开子集的有限交，这是 $\scr{U}_a$ 的一个稠开子集”——$\scr{U}'\subset\scr{U}_{a_i}$ 本无包含之碍，误在\trDense：交 $\scr{U}'$ 未必于单独某个 $\scr{U}_{a_i}$ 中\trDense，应取共同邻域 $\bigcap_i\scr{U}_{a_i}$ 言其稠且开（存疑照排；译注 2026-09-10 改写）。
\item[p.~546] 定理 A.3.10 证明 $G$ 的第二支印作 $g(x)+\alpha(v)v$ $\to$ $\alpha$ 为 $x$ 之函数、$v$ 为参数变量，$\alpha(v)$ 无定义，应作 $g(x)+\alpha(x)v$（径改）。
\item[p.~548] “Cartans' Formula” $\to$ “Cartan's Formula”（径改）。
\item[p.~549] 引理 A.4.2 证明印作 $\scr{L}_{X}\widetilde{Y}$（$=[\widetilde{Y},X]$）$\to$ 按本书括号约定 $[X,Y]f=X(Yf)-Y(Xf)$ 应作 $[X,\widetilde{Y}]$（原式与 $D(fY)=\dot{f}Y+fDY$ 及 p.~551 的 $A(t)=-(dX)$ 反号相抵）（径改）。
\item[p.~550] 变分方程处的 $\mathbf{C}^m$ $\to$ $\mathbf{R}^m$（径改）。
\item[p.~551] $\widetilde{X}$ 的流 $\widetilde{\varphi}^{t}(u,x)=(u+t,\varphi^{t}(x))$ 及下文两处切流等式 $\to$ 应含初始时间贡献：$(u+t,\,\varphi^{u+t}\circ(\varphi^{u})^{-1}(x))$（反例 $X_{t}=t\partial_{x}$），切流基点应作 $x(0)$（存疑照排）。
\item[p.~552] $(x\cdot f)Y$ $\to$ $(X\cdot f)Y$（径改）。
\item[p.~552] $D_{\dot\gamma}\gamma=0$ $\to$ $D_{\dot\gamma}\dot\gamma=0$（径改）。

\subsection*{附录~B}
\item[p.~556] $0\to D_{j-1}\xrightarrow{\ \partial\ }D_{j-1}\to0$ 首项 $\to$ $D_j$（径改）。
\item[p.~556] $\partial_{i+j+1}\colon C_{i+1}\otimes D_{j}\longrightarrow C_{i}\otimes D_{j}$ $\to$ 其靶应为整个 $(C\otimes D)_{i+j}$，漏排第二个\trDirectSum{}项 $C_{i+1}\otimes D_{j-1}$（下文所求边缘的第二分量恰非零地落在该分量中）（存疑照排）。
\item[p.~556] $\partial y_a^{i}$ $\to$ $\partial y_a^{j}$（径改）。
\item[p.~557] $E_{j-1}$（无撇） $\to$ $E'_{j-1}$（径改）。

\subsection*{附录~C}
\item[p.~566] 命题 C.2.7 证明收尾对 $E$ 中任意序列 $(x_k)$ 由 $L(x_k)\to y$ 断言其为 \trCauchy{}序列照排——估计 $\|x\|\le C\|L(x)\|$ 仅对 $E_0$ 中的 $x$ 建立，$\operatorname{Ker}L\ne\{0\}$ 时推理失效（取 $x_k=kv$，$v\in\operatorname{Ker}L$）$\to$ 注载修补：换用其在闭补 $E_0$ 中的代表元（命题 8.7.4 无界情形所用的正是此手法），结论不受影响（存疑照排）。
\item[p.~572] 一处 $\varphi$ 的全空间印作粗正体 $\mathbf{C}^\infty_0$ $\to$ 花体 $\cC^\infty_0$（径改）。
\item[p.~573] “方程 $Lu=0$ 的\trFundamentalSolution” $\to$ 按定义式应为 $Lu=\delta$（径改）。
\item[p.~573] 定理 C.3.1 称方程 $Lu=f$ 拥有唯一的紧\trSupport\trWeakSolution{} $u=E\star f$ 照排——存在性可失效（$L=\frac{d}{dx}$、$E$ 为 Heaviside 函数、$f=\delta$：$E\star f$ 非紧支，且紧支解不存在），证明恰只证了条件唯一性那一半（存疑照排）。
\item[p.~576] 两处 display 被积函数 $u(\tau)/d\tau$ $\to$ $du/d\tau$（径改）；同注并案：次式内层积分 $\int_{t}^{t_{1}}$ 有向，$t>t_{1}$ 时其值为负、对全方形积分后中项恒为零（$u(t)=t-\frac{1}{2}$ 即反例），积分限应改作 $\min(t,t_{1})$ 至 $\max(t,t_{1})$ 再放大（存疑照排）。
\item[p.~577/579] 定理 C.4.6 与定理 C.4.10 照排并案——仅设 $U$ 为\trBounded\trOpen{}子集时\trCompact{}性断言不成立（两两不交、半径趋零的开球之并上取规范化的分段常值函数列，无强\trConvergence{}子列），应补边界\trRegularity{}假设或相当的延拓条件（存疑照排）。
\item[p.~578] 定理 C.4.9 的 $|\alpha|\le k$ $\to$ 疑应限于满足 $k-|\alpha|>n/p$ 的那些 $\alpha$（\trSobolev\trEmbedding{}的严格范围；否则 $u(x)=x^{3/4}\in W^{1,2}(0,1)$ 而 $u'\notin L^{\infty}$ 即成反例）；末句 $m=[k-n/p]$ 于 $k-n/p$ 恰为整数时须另附加\trRegularity{}条件（存疑照排）。
\item[p.~580] \trCauchy{}核计算两处微元印作 $d\rho\,d\theta$ $\to$ 按积分限（外层 $0$ 至 $+\infty$ 配 $\rho$、内层 $0$ 至 $2\pi$ 配 $\theta$）先后次序应反，应作 $d\theta\,d\rho$（存疑照排）。
\item[p.~581] $\chi_{r}\star\chi_{r}$ 显式式首项 $2\arccos\frac{\|x\|}{2r}$ $\to$ 漏因子 $r^{2}$（以 $x=0$ 核验，交叠面积应为 $\pi r^{2}$ 而原式给 $\pi$），应作 $2r^{2}\arccos\frac{\|x\|}{2r}$（径改）。
\item[p.~582] 韦尔引理证明称 $u\star\beta_{\delta}$ 的\trSupport{}含于 $U_{\delta}=\{x\in U\mid B(x,\delta)\subset U\}$ 照排——$U_{\delta}$ 系\trConvolution{}的定义域而非\trSupport{}所在（取 $U$ 为\trDisk{}、$u\equiv1$ 即反例），应述“$u\star\beta_{\delta}$ 在 $U_{\delta}$ 上有定义且属 $\cC^{\infty}$ 类”（存疑照排）。
\item[p.~583] 平均值不等式证明中三处 $B_\rho$ $\to$ $\partial B_\rho$（径改）。
\item[p.~583] “$K$ 是\trHarmonic”句与上文 $\Delta K=\delta$ 的关系存疑，疑指 $K$ 在 $\CC\setminus\{0\}$ 上\trHarmonic（存疑照排）。
\item[p.~584] $\frac{\partial}{\partial s}\Delta K=\frac{\partial}{\partial s}\,\delta$、$\frac{\partial}{\partial t}\,\delta$ $\to$ $\frac{\partial}{\partial x}$、$\frac{\partial}{\partial y}$（径改）。
\item[p.~584] $\|\operatorname{grad}(K_{j}\star f)\|_{L^{p}}\le C(p)\,\|\Delta f\|_{L^{p}}$ 右端与左端的缩放尺度不\trCompatible（缩放反例可直见）$\to$ 按左端（零阶\trSingular{}积分）右端疑应作 $\|f\|_{L^{p}}$（存疑照排）。

\end{description}
